\setcounter{section}{12}

  \zwischenueberschrift{Bundel vektor}

Misalkan $M$ suatu
\definitionsverweis {manifold diferensiabel}{}{}
dengan
\definitionsverweis {bundel tangen}{}{}
\maabb {p} {TM} {M
} {.}
Untuk suatu domain peta

\mavergleichskette
{\vergleichskette
{ U
}
{ \subseteq }{ M
}
{  }{
}
{    }{
}
{   }{
}
}
{}{}{}
dan suatu peta
\maabbdisp {\alpha} { U } { V \subseteq \R^n
} {}
berdasarkan
[[Differenzierbare Mannigfaltigkeit/Abbildung/Tangentialabbildung/Eigenschaften/Fakt|Lema 10.8]]
terdapat suatu
\definitionsverweis {homeomorfisme}{}{}
\maabbdisp {T(\alpha)} {TM\vert_U \cong TU} { TV \cong V \times \R^n
} {,}
yang linear pada komponen kedua untuk setiap titik dasar

\mavergleichskette
{\vergleichskette
{ P
}
{ \in }{ U
}
{  }{
}
{    }{
}
{   }{
}
}
{}{}{}
. Untuk peta lain
\maabbdisp {\alpha'} {U'} {V'
} {}
pemetaan transisinya
\maabbdisp {} {TV \vert_{\alpha(U \cap U') }  \cong \alpha(U \cap U') \times \R^n } { TV' \vert_{\alpha'(U \cap U') }  \cong \alpha'(U \cap U') \times \R^n
} {}
linear pada serat-seratnya. Melalui homeomorfisme

\mavergleichskette
{\vergleichskette
{ U
}
{ \cong }{ V
}
{  }{
}
{    }{
}
{   }{
}
}
{}{}{}
kita juga langsung memperoleh homeomorfisme

\mavergleichskettedisp
{\vergleichskette
{ TM \vert_U \cong TU
}
{ \cong} { U \times \R^n
}
{ } {
}
{ } {
}
{ } {
}
}
{}{}{}
di atas $U$ beserta pemetaan transisi terkait
\maabbdisp {} {(U \cap U') \times \R^n } { (U \cap U') \times \R^n
} {,}
yang linear terhadap $\R^n$.

Bundel tangen merupakan contoh penting dari bundel vektor.

\inputdefinition
{}
{

Misalkan $X$ suatu
\definitionsverweis {ruang topologis}{}{}
dan

\mavergleichskette
{\vergleichskette
{r
}
{ \in }{\N
}
{  }{
}
{    }{
}
{   }{
}
}
{}{}{.}
Sebuah
\definitionswort {bundel vektor riil}{}
$V$ ber-
\definitionswort {rank}{}
$r$ adalah suatu ruang topologis yang dilengkapi dengan
\definitionsverweis {pemetaan kontinu}{}{}
\maabb {p} { V } { X
} {}
sedemikian sehingga setiap
\definitionsverweis {serat}{}{}
$p^{-1}(x)$ merupakan
\definitionsverweis {ruang vektor riil}{}{}
yang
\definitionsverweis {berdimensi}{}{}
$r$ dan terdapat suatu
\definitionsverweis {penutup terbuka}{}{}

\mavergleichskette
{\vergleichskette
{X
}
{ = }{ \bigcup_{i \in I} U_i
}
{  }{
}
{    }{
}
{   }{
}
}
{}{}{}
beserta
\definitionsverweis {homeomorfisme-homeomorfisme}{}{}
\maabbdisp {\varphi_i} {p^{-1} \left(  U_i  \right) } { U_i \times \R^r
} {}
di atas $U_i$ yang pada setiap serat menginduksi suatu
\definitionsverweis {isomorfisme linear}{}{}
\maabbdisp {\left(  \varphi_i  \right)_x} {p^{-1} (x) } { \R^r
} {}
.

}

Dalam hal ini, $V$ juga disebut \stichwort {ruang total} {} dan $X$ disebut \stichwort {ruang dasar} {} dari bundel vektor tersebut. Untuk serat sering pula ditulis

\mavergleichskette
{\vergleichskette
{ V_x
}
{ = }{ p^{-1} (x)
}
{  }{
}
{    }{
}
{   }{
}
}
{}{}{.}

Dalam homeomorfisme
\maabb {} { p^{-1}(U) } { U \times \R^r
} {}
sisi kanan dilengkapi dengan
\definitionsverweis {topologi produk}{}{}
dan $\R^r$ dengan
\definitionsverweis {topologi Euklides}{}{}
alaminya, sedangkan
\mathl{p^{-1}(U)}{} dilengkapi dengan topologi yang diinduksi dari $V$. Dengan demikian, semua serat $V_x$ menyandang topologi alami suatu ruang vektor riil berdimensi hingga. Yang dimaksud dengan homeomorfisme di atas $U$ adalah bahwa diagram

\mathdisp {\begin{matrix} p^{-1}(U)   & \stackrel{  }{\longrightarrow} &  U \times \R^r   & \\ &  \searrow & \downarrow & \\ & &  U   & \!\!\!\!\! \!\!\!  \\ \end{matrix}} {  }

komutatif. Produk
\mathl{X \times \R^r}{} merupakan suatu bundel vektor yang disebut \stichwort {bundel vektor trivial} {.} Suatu homeomorfisme linear
\maabb {} { p^{-1}(U) } { U \times \R^r
} {}
disebut
\zusatzklammer {lokal} {} {}
\stichwort {trivialisasi} {.} Pembatasan suatu bundel vektor pada $U_i$ bersifat trivial; jadi, secara lokal setiap bundel vektor bersifat trivial. Bundel vektor ber-rank $1$ juga disebut \stichwort {bundel garis} {.}

\inputbemerkung
{}
{

Suatu
\definitionsverweis {bundel vektor riil}{}{}
\maabb {p} { V } { X
} {}
di atas
\definitionsverweis {ruang topologis}{}{}
$X$ juga dapat didefinisikan hanya dengan mengacu pada suatu penutup terbuka

\mavergleichskettedisp
{\vergleichskette
{ X
}
{ =} { \bigcup_{i \in I} U_i
}
{ } {
}
{ } {
}
{ } {
}
}
{}{}{}
dan trivialisasi-trivialisasi
\maabbdisp {\varphi_i} {p^{-1} \left(  U_i  \right) } { U_i \times \R^r
} {}
dengan mensyaratkan agar pemetaan transisi
\zusatzklammer {yang diperoleh di atas $p^{-1} (U_i \cap U_j)$} {} {}
\maabbdisp {\varphi_j \circ \varphi_i^{-1}} {U_i \cap U_j \times \R^r } { U_i \cap U_j \times \R^r
} {}
linear pada komponen ruang vektornya. Hal ini menghasilkan struktur ruang vektor pada setiap serat, dengan menggunakan sebarang trivialisasi.

}

__NOEDITSECTION__

\inputfaktbeweis
{Topologischer Raum/Reelles Vektorbündel/Einschränkung/Fakt}
{Lema}
{}
{

\faktsituation {Misalkan
\maabb {p} { V } { X
} {}
suatu
\definitionsverweis {bundel vektor riil}{}{}
di atas
\definitionsverweis {ruang topologis}{}{}
$X$.}
\faktfolgerung {Maka untuk setiap
\definitionsverweis {himpunan terbuka}{}{}

\mavergleichskette
{\vergleichskette
{ W
}
{ \subseteq }{ X
}
{  }{
}
{    }{
}
{   }{
}
}
{}{}{}
\definitionsverweis {pembatasan}{}{}
\maabbdisp {} {p^{-1} (W)} {W
} {}
juga merupakan suatu bundel vektor.}
\faktzusatz {}
\faktzusatz {}

}
{

Hal ini jelas; kita cukup membatasi homeomorfisme-homeomorfisme yang linear pada serat
\maabbdisp {\varphi_i} {p^{-1} \left(  U_i  \right) } { U_i \times \R^r
} {}
menjadi
\maabbdisp {\varphi_i \vert_{W \cap U_i }} {p^{-1} \left(  W \cap U_i  \right) } { \left(  W \cap U_i  \right)  \times \R^r
} {}
.

}

\inputdefinition
{}
{

Misalkan
\mathkor {} {E} {dan} {F} {}
\definitionsverweis {bundel-bundel vektor riil}{}{}
di atas suatu
\definitionsverweis {ruang topologis}{}{}
$X$. Suatu
\definitionswort {homomorfisme bundel vektor}{}
\maabb {\varphi} { E } { F
} {}
adalah
\definitionsverweis {pemetaan kontinu}{}{}
di atas $X$ sedemikian sehingga untuk setiap titik

\mavergleichskette
{\vergleichskette
{ x
}
{ \in }{ X
}
{  }{
}
{    }{
}
{   }{
}
}
{}{}{}
pemetaan terinduksi
\maabbdisp {\varphi_x} { E_x} {F_x
} {}
bersifat
$\R$-\definitionsverweis {linear}{}{}
.

}

Pengamatan berikut memungkinkan konstruksi berbagai macam bundel vektor.

\inputbemerkung
{}
{

Misalkan $X$ suatu
\definitionsverweis {ruang topologis}{}{}
dan misalkan

\mavergleichskette
{\vergleichskette
{ f_1 , \ldots , f_n
}
{ \in }{ C^0(X,\R)
}
{  }{
}
{    }{
}
{   }{
}
}
{}{}{}
merupakan
\definitionsverweis {fungsi-fungsi kontinu}{}{}
bernilai riil pada $X$. Fungsi-fungsi ini mendefinisikan suatu pemetaan
\maabbeledisp {} {X \times \R^n } { X \times \R
} {(P,t_1 , \ldots , t_n) } { (P, \sum_{i= 1}^n f_i (P) t_i)
} {,}
antara
\definitionsverweis {bundel vektor}{}{}
trivial ber-rank $n$ di atas $X$ dan bundel vektor trivial ber-rank $1$ di atas $X$. Pemetaan ini kontinu, dan di atas setiap titik

\mavergleichskette
{\vergleichskette
{ P
}
{ \in }{ X
}
{  }{
}
{    }{
}
{   }{
}
}
{}{}{}
terdapat
\definitionsverweis {bentuk linear}{}{}
\maabb {} {\R^n } { \R
} {}
dengan matriks baris
\mathl{\left( f_1(P)  , \,   \ldots  , \,  f_n(P)                 \right)}{}; oleh karena itu, pemetaan ini merupakan suatu
\definitionsverweis {homomorfisme}{}{}
bundel vektor. Melalui kernel tersebut, pada setiap serat $\R^n$ diperoleh suatu subruang vektor berdimensi $n-1$, kecuali apabila semua fungsi $f_i$ lenyap secara serentak di titik $P$. Misalkan

\mavergleichskette
{\vergleichskette
{ D(f_i)
}
{ = }{ \left\{ P \in X  \mid  f_i(P) \neq 0        \right\}
}
{  }{
}
{    }{
}
{   }{
}
}
{}{}{,}
maka

\mavergleichskette
{\vergleichskette
{ U
}
{ = }{ \bigcup_{i = 1}^n D(f_i)
}
{ \subseteq }{ X
}
{    }{
}
{   }{
}
}
{}{}{}
adalah himpunan bagian terbuka tempat sekurang-kurangnya satu $f_i$ tidak lenyap, sehingga kernel-kernel tersebut selalu berdimensi $n-1$. \stichwort {Bundel kernel} {} $E$ yang terkait dengan fungsi-fungsi $f_i$ adalah bundel vektor di atas $U$ yang ditentukan serat demi serat oleh kernel-kernel itu, yaitu

\mavergleichskettedisp
{\vergleichskette
{ E
}
{ =} { \left\{  (P; t_1 , \ldots , t_n)  \mid   P \in U , \,  \sum_{i= 1}^n f_i(P)t_i = 0       \right\}
}
{ \subseteq} { U \times \R^n
}
{ } {
}
{ } {
}
}
{}{}{.}
Berdasarkan
[[Topologischer Raum/Triviale Vektorbündel/Homomorphismus/Kernbündel/Trivialisierungen/Aufgabe|Soal 12.26]]
terdapat trivialisasi-trivialisasi di atas $D(f_i)$, dan objek ini memang merupakan suatu bundel vektor.

}

\inputbeispiel{}
{

Di atas $\R^2$, kita meninjau pemetaan
\maabbeledisp {} {\R^2 \times \R^2} { \R^2 \times \R
} {(x,y;s,t)} { (x,y;xs+yt)
} {,}
sebagai suatu
\definitionsverweis {homomorfisme bundel vektor}{}{}
dalam pengertian
[[Topologischer Raum/Triviale Vektorbündel/Homomorphismus/Kernbündel/Bemerkung|Catatan 12.5]].
Satu-satunya nol bersama kedua fungsi koordinat $x,y$ adalah titik asal $(0,0)$, sehingga

\mavergleichskette
{\vergleichskette
{ U
}
{ = }{ \R^2 \setminus \{ (0,0) \}
}
{  }{
}
{    }{
}
{   }{
}
}
{}{}{}
dan bundel kernel terkait adalah

\mavergleichskettedisp
{\vergleichskette
{ E
}
{ =} { \left\{  (P; s ,t)   \mid   P \in U  , \,  xs + yt = 0       \right\}
}
{ \subseteq} { U \times \R^2
}
{ } {
}
{ } {
}
}
{}{}{.}
Bundel ini terdiri atas suatu keluarga garis di dalam suatu bidang, dengan garis-garis tersebut bergerak bersama titik dasar

\mavergleichskette
{\vergleichskette
{ P
}
{ = }{ (x,y)
}
{  }{
}
{    }{
}
{   }{
}
}
{}{}{}
.

}

\inputbeispiel{}
{

Misalkan

\mavergleichskette
{\vergleichskette
{ W
}
{ \subseteq }{ \R^n
}
{  }{
}
{    }{
}
{   }{
}
}
{}{}{}
\definitionsverweis {terbuka}{}{,}
\maabb {h} { W } { \R
} {}
suatu fungsi yang dua kali
\definitionsverweis {terdiferensialkan secara kontinu}{}{}
dan

\mavergleichskette
{\vergleichskette
{ Y
}
{ = }{ h^{-1}(c)
}
{  }{
}
{    }{
}
{   }{
}
}
{}{}{}
merupakan
\definitionsverweis {serat}{}{}
di atas

\mavergleichskette
{\vergleichskette
{ c
}
{ \in }{ \R
}
{  }{
}
{    }{
}
{   }{
}
}
{}{}{,}
dengan $h$
\definitionsverweis {reguler}{}{}
di setiap titik dari $Y$. \definitionsverweis {Ruang singgung}{}{}
$T_PY$ pada

\mavergleichskette
{\vergleichskette
{ P
}
{ \in }{ Y
}
{  }{
}
{    }{
}
{   }{
}
}
{}{}{}
adalah kernel
\definitionsverweis {diferensial total}{}{}
$\left(D h \right)_{ P}$, yang diberikan oleh turunan-turunan parsial ${ \frac{   \partial   h   }{  \partial   x_1   } } ( P) , \ldots , { \frac{   \partial   h   }{  \partial   x_n   } } ( P)$ dan karenanya merupakan subruang vektor berdimensi $(n-1)$ dari $\R^n$; lihat juga
[[Abgeschlossene Untermannigfaltigkeit/Punktweise/Tangentialraum als Unterraum/Fakt|Teorema 10.3]].
Kita meninjau pemetaan
\maabbeledisp {} {Y \times \R^n } { Y \times \R
} {(P;t_1 , \ldots , t_n) } { (P; \sum_{i= 1}^n { \frac{   \partial   h   }{  \partial   x_i   } } ( P)  t_i )
} {}
dan bundel kernel terkait

\mavergleichskettedisp
{\vergleichskette
{ E
}
{ =} { \left\{  (P; t_1 , \ldots , t_n)  \mid   P \in Y , \,  \sum_{i= 1}^n { \frac{   \partial   h   }{  \partial   x_i   } } ( P) t_i = 0       \right\}
}
{ \subseteq} { Y \times \R^n
}
{ } {
}
{ } {
}
}
{}{}{.}
dalam pengertian
[[Topologischer Raum/Triviale Vektorbündel/Homomorphismus/Kernbündel/Bemerkung|Catatan 12.5]].
Bundel ini secara langsung berimpit, serat demi serat, dengan ruang-ruang singgung pada $Y$. Namun, bundel ini memang merupakan bundel tangen, sebagaimana diperoleh dengan membandingkannya dengan
[[Abgeschlossene Untermannigfaltigkeit/Tangentialbündel/Abgeschlossen/Aufgabe|Soal 10.22]].

}

\inputdefinition
{}
{

Misalkan
\mathkor {} {E} {dan} {F} {}
\definitionsverweis {bundel-bundel vektor riil}{}{}
di atas suatu
\definitionsverweis {ruang topologis}{}{}
$X$. Suatu
\definitionsverweis {homomorfisme bundel vektor}{}{}
\maabb {\varphi} { E } { F
} {}
disebut
\definitionsverweis {isomorfisme}{}{,}
apabila terdapat suatu homomorfisme
\maabb {\psi} { F } { E
} {}
yang, ketika dikomposisikan dengan $\varphi$
\zusatzklammer {dalam kedua urutan} {} {}, menghasilkan identitas.

}

\inputdefinition
{}
{

Misalkan
\mathkor {} {X} {dan} {Y} {}
\definitionsverweis {ruang-ruang topologis}{}{}
dan misalkan
\maabbdisp {p} { Y } { X
} {}
suatu
\definitionsverweis {pemetaan kontinu}{}{.}
Yang dimaksud dengan
\definitionswort {penampang kontinu}{}
dari $p$ adalah suatu pemetaan kontinu
\maabb {s} { X } { Y
} {}
dengan

\mavergleichskettedisp
{\vergleichskette
{  p \circ s
}
{ =} {
\operatorname{Id}_{  X  }
}
{ } {
}
{ } {
}
{ } {
}
}
{}{}{.}

}
Sebagai contoh, $Y$ dapat dipandang sebagai suatu bundel vektor di atas $X$. Suatu penampang hanya mungkin ada apabila $p$ surjektif, dan ini selalu berlaku bagi suatu bundel vektor. Kadang-kadang suatu penampang diidentifikasi dengan citranya; hal ini tidak menimbulkan masalah karena suatu penampang selalu injektif. Peran khusus dimainkan oleh \stichwort {penampang nol} {,} yang memetakan setiap titik dasar $P$ ke titik nol dalam ruang vektor $V_P$.
Untuk
\definitionsverweis {bundel tangen}{}{}
penampang-penampang kontinu tepat merupakan
\definitionsverweis {medan-medan vektor}{}{}
kontinu. Untuk bundel vektor trivial
\maabbdisp {} {X \times \R^n } { X
} {}
terdapat penampang-penampang konstan

\mathbed {e_i} {}
 {i = 1 , \ldots , n} {}
 {} {} {} {,}
yang memetakan setiap titik ke vektor

\mavergleichskette
{\vergleichskette
{ e_i
}
{ \in }{ \R^n
}
{  }{
}
{    }{
}
{   }{
}
}
{}{}{}
. Keluarga penampang ini membentuk suatu basis pada serat di atas setiap titik.

Misalkan
\maabb {p} { V } { X
} {}
suatu
\definitionsverweis {bundel vektor riil}{}{}
ber-
\definitionsverweis {rank}{}{}
$r$ di atas suatu
\definitionsverweis {ruang topologis}{}{}
$X$, dan misalkan

\mavergleichskette
{\vergleichskette
{ U
}
{ \subseteq }{ X
}
{  }{
}
{    }{
}
{   }{
}
}
{}{}{}
suatu
\definitionsverweis {himpunan bagian terbuka}{}{.}
Suatu keluarga
\definitionsverweis {penampang}{}{}
kontinu
\maabbdisp {s_1 , \ldots , s_r} { U} {V  \vert_U
} {}
disebut keluarga
\definitionswort {penampang basis}{,}
apabila untuk setiap titik

\mavergleichskette
{\vergleichskette
{ P
}
{ \in }{ U
}
{  }{
}
{    }{
}
{   }{
}
}
{}{}{}
vektor-vektor

\mavergleichskette
{\vergleichskette
{ s_1(P)  , \ldots , s_r(P)
}
{ \in }{ V_P
}
{  }{
}
{    }{
}
{   }{
}
}
{}{}{}
membentuk suatu
\definitionsverweis {basis}{}{}
dari $V_P$.

Keberadaan penampang-penampang basis
\zusatzklammer {kontinu} {} {}
pada suatu himpunan terbuka $U$ ekuivalen dengan keberadaan suatu trivialisasi
\zusatzklammer {kontinu} {} {}
dari bundel vektor di atas $U$; lihat
[[Vektorbündel/Offene Menge/Stetige Basisschnitte/Trivialisierung/Aufgabe|Soal 12.14]].

  \zwischenueberschrift{Data perekatan}

Gagasan untuk merekatkan data lokal guna mendeskripsikan objek global telah kita jumpai dalam
[[Topologische Mannigfaltigkeit/Rekonstruktion aus Karten/Aufgabe|Soal 8.15]].

\inputdefinition
{}
{

Yang dimaksud dengan suatu
\definitionswort {data perekatan}{}
untuk
\definitionsverweis {bundel vektor riil}{}{}
ber-rank $r$ di atas suatu
\definitionsverweis {ruang topologis}{}{}
$X$ adalah kumpulan data berikut.
\aufzaehlungvier{Suatu
\definitionsverweis {penutup terbuka}{}{}

\mavergleichskettedisp
{\vergleichskette
{X
}
{ =} { \bigcup_{i \in I} U_i
}
{ } {
}
{ } {
}
{ } {
}
}
{}{}{}
}{Suatu keluarga

\mathbed {E_i \rightarrow U_i} {}
 {i \in I} {}
 {} {} {} {,}
bundel vektor riil ber-rank $r$.
}{Untuk setiap pasangan
\mathl{(i,j)}{} suatu
\definitionsverweis {isomorfisme bundel vektor}{}{}
\maabbdisp {\varphi_{ji}} { E_i {{|}}_{U_i \cap U_j }
} { E_{j} {{|}}_{U_i \cap U_j }
} {}
di atas
\mathl{U_i \cap U_j}{.}
}{Untuk indeks-indeks

\mavergleichskette
{\vergleichskette
{ i,j,k
}
{ \in }{ I
}
{  }{
}
{    }{
}
{   }{
}
}
{}{}{}
berlaku
\definitionswort {syarat koklus}{}

\mavergleichskettedisp
{\vergleichskette
{ \varphi_{kj} \circ \varphi_{ji}
}
{ =} { \varphi_{ki}
}
{ } {
}
{ } {
}
{ } {
}
}
{}{}{}
sebagai pemetaan dari
\mathl{E_{i} {{|}}_{U_i \cap U_j \cap U_k }}{} ke
\mathl{E_{k}  {{|}}_{U_i \cap U_j \cap U_k }}{}.
}

}

\inputbemerkung
{}
{

Secara khas, dalam
[[Topologisches reelles Vektorbündel/Verklebungsdatum/Definition|Definisi 12.10]]
bundel-bundel vektor pada (2) adalah bundel-bundel vektor trivial di atas $U_i$, yakni

\mavergleichskette
{\vergleichskette
{E_i
}
{ = }{ \R^r \times U_i
}
{  }{
}
{    }{
}
{   }{
}
}
{}{}{.}
Isomorfisme-isomorfisme pada (3) kemudian hanya diberikan oleh pemetaan linear bijektif
\maabb {\varphi_{ji}} { \R^r} { \R^r
} {}
yang bergantung secara kontinu pada titik dasar dari
\mathl{U_i \cap U_j}{}. Pemetaan-pemetaan ini dapat dideskripsikan secara ringkas oleh pemetaan kontinu
\maabbdisp {\varphi_{ji}} {U_i \cap U_j } { \operatorname{GL}_{  r   } \! \left(  \R  \right)
} {}
ke dalam
\definitionsverweis {grup linear umum}{}{}
. Jadi, setiap titik dasar dipetakan secara kontinu ke suatu objek
\definitionsverweis {invertibel}{}{}
berupa $r \times r$-
\definitionsverweis {matriks}{}{,}
di mana kontinuitas berarti bahwa semua entri matriks merupakan fungsi kontinu. Ini disebut \stichwort {deskripsi matriks} {} dari bundel tersebut. Syarat koklus tetap berlaku.

}

__NOEDITSECTION__
\inputfaktbeweis
{Topologisches reelles Vektorbündel/Verklebungsdatum/Existenz/Fakt}
{Lema}
{}
{

\faktsituation {Misalkan diberikan suatu
\definitionsverweis {data perekatan}{}{}

\mathbed {E_i} {}
 {i \in I} {}
 {} {} {} {,}
di atas suatu
\definitionsverweis {ruang topologis}{}{}

\mavergleichskettedisp
{\vergleichskette
{X
}
{ =} { \bigcup_{i \in I} U_i
}
{ } {
}
{ } {
}
{ } {
}
}
{}{}{}
.}
\faktfolgerung {Maka terdapat suatu bundel vektor riil yang ditentukan secara tunggal
\maabb {} { E } { X
} {}
beserta
\definitionsverweis {isomorfisme-isomorfisme}{}{}
\maabb {\psi_i} {E_i } { E \vert_{U_i }
} {}
sedemikian sehingga

\mavergleichskettedisp
{\vergleichskette
{ \psi_i \vert_{ E_i {{ |}}_{ U_{i} \cap U_j } }
}
{ =} { \psi_j \vert_{ E_j {{ |}}_{ U_{i} \cap U_j } }  \circ \varphi_{ji}
}
{ } {
}
{ } {
}
{ } {
}
}
{}{}{}
berlaku.}
\faktzusatz {}
\faktzusatz {}

 }
{ Lihat [[Topologisches reelles Vektorbündel/Verklebungsdatum/Existenz/Fakt/Beweis/Aufgabe|Soal 12.12]]. }

\bild{ \begin{center}
\includegraphics[width=5.5cm]{ \bildeinlesunggif {Fiddler_crab_mobius_strip} {gif} }
\end{center}
\bildtext {Animasi seekor kepiting biola yang mengelilingi pita Möbius dan kembali dengan orientasi tercermin. Edisi PDF menampilkan bingkai pertama yang statis; animasi asli dipertahankan untuk edisi HTML dan unduhan.} }

\bildlizenz { Fiddler_crab_mobius_strip.gif } {} {Hamishtodd1} {Commons} {CC-by-sa 4.0} {}

\inputbeispiel{}
{

Pada lingkaran satu dimensi

\mavergleichskettedisp
{\vergleichskette
{S^1
}
{ =} { \left\{ (x,y) \in \R^2  \mid   x^2+y^2 = 1        \right\}
}
{ } {
}
{ } {
}
{ } {
}
}
{}{}{}
kita meninjau
\definitionsverweis {penutup terbuka}{}{}

\mavergleichskette
{\vergleichskette
{S^1
}
{ = }{ U \cup V
}
{  }{
}
{    }{
}
{   }{
}
}
{}{}{}
dengan

\mavergleichskette
{\vergleichskette
{U
}
{ = }{S^1 \setminus \{(0,1)\}
}
{  }{
}
{    }{
}
{   }{
}
}
{}{}{}
dan

\mavergleichskette
{\vergleichskette
{V
}
{ = }{S^1 \setminus \{(0,-1)\}
}
{  }{
}
{    }{
}
{   }{
}
}
{}{}{.}
Di atasnya kita mendeskripsikan suatu
\definitionsverweis {data perekatan}{}{}
untuk bundel vektor riil ber-rank $1$. Kedua himpunan terbuka tersebut
\definitionsverweis {homeomorfik}{}{}
dengan garis riil. Kita mempunyai

\mavergleichskettedisp
{\vergleichskette
{U \cap V
}
{ =} { S^1 \setminus \{(0,1), (0,-1)  \}
}
{ =} { \left\{ (x,y) \in  S^1   \mid   x \neq 0        \right\}
}
{ } {
}
{ } {
}
}
{}{}{,}
dan himpunan ini tidak
\definitionsverweis {terhubung}{}{,}
melainkan homeomorfik dengan dua setengah garis riil terbuka yang saling lepas
\zusatzklammer {atau, ekuivalen, dengan garis-garis} {} {.}
Kita tetapkan

\mavergleichskette
{\vergleichskette
{L
}
{ = }{ U \times \R
}
{  }{
}
{    }{
}
{   }{
}
}
{}{}{}
dan

\mavergleichskette
{\vergleichskette
{M
}
{ = }{ V \times \R
}
{  }{
}
{    }{
}
{   }{
}
}
{}{}{.}
Kita mendefinisikan suatu
\definitionsverweis {isomorfisme}{}{}
\maabbdisp {\varphi} {L \vert_{U \cap V} } { M\vert_{U \cap V}
} {}
melalui

\mavergleichskettedisp
{\vergleichskette
{ \varphi(x,y,t)
}
{ \defeq} { \begin{cases}  (x,y,t) \text{ untuk } x > 0 \, , \\  (x,y,-t) \text{ untuk } x < 0  \,  \end{cases}
}
{ } {
}
{ } {
}
{ } {
}
}
{}{}{}
. Perhatikan bahwa $\varphi$ kontinu karena kedua ekspresi fungsional itu berlaku pada himpunan bagian terbuka yang saling lepas. Pada separuh yang satu pemetaannya adalah identitas, sedangkan pada separuh yang lain pemetaannya membalik. Dalam pengertian
[[Verklebungsdatum/Vektorbündel/Trivialisierungen/Bemerkung|Catatan 12.11]]
kita memperoleh deskripsi matriks
\zusatzklammer {konstan} {} {}
yang kontinu

\mavergleichskettedisp
{\vergleichskette
{ \psi (x,y)
}
{ \defeq} { \begin{cases}  (1) \text{ untuk } x > 0 \, , \\  (-1) \text{ untuk } x < 0  \, , \end{cases}
}
{ } {
}
{ } {
}
{ } {
}
}
{}{}{}
pada
\mathl{U \cap V}{}. Karena hanya terdapat dua himpunan terbuka, syarat koklus terpenuhi secara otomatis. Berdasarkan
[[Topologisches reelles Vektorbündel/Verklebungsdatum/Existenz/Fakt|Lema 12.12]]
data perekatan ini mendefinisikan suatu bundel vektor riil ber-rank $1$ pada lingkaran tersebut, yang disebut \stichwort {pita Möbius} {}.

}

 [[Kategori:Latexseite]]
